\chapter{Diseño del Sistema y Decisiones Técnicas}
\label{cap:disenoSistema}

En este capítulo se detallan las decisiones técnicas adoptadas para la implementación de ClimbIt
así como la justificación de dichas elecciones frente a las alternativas actuales del mercado. 
También se detallarán las arquitecturas del frontend y del backend y los patrones de diseño 
utilizados.

\section{Decisiones de Diseño y Stack Tecnológico}

Para ClimbIt surgieron varias opciones de diseño para el sistema, pero al final optamos por 
una separación de responsabilidades para ceder la lógica de negocio y la persistencia de los datos 
al backend y un frontend SPA (\textit{Single Page Application}) simple y de fácil navegación 
que consuma esos datos a través de una API. 
Desde el comienzo sabíamos que el alcance de esta aplicación era notablemente amplio, por lo que 
la decisión sobre qué arquitectura utilizar fue previamente meditada. Valoramos tres opciones de 
diseño para el backend: una arquitectura monolítica, una arquitectura hexagonal y 
una arquitectura de microservicios. 

Para la elección tuvimos en cuenta tres factores principales: la escalabilidad, la mantenibilidad y la complejidad 
de desarrollo. Una arquitectura monolítica podría haber sido más sencilla de implementar inicialmente, pero a 
medida que el proyecto creciese, podría habernos llevado a un código difícil de mantener y escalar. Por otro lado, 
una arquitectura de microservicios nos habría ofrecido grandes ventajas en cuanto a la escalabilidad, pero la curva de 
aprendizaje y la dificultad que añadía al desarrollo y la necesidad de gestionar múltiples servicios nos 
hizo descartarla, por el momento, para este proyecto. Al final decidimos que la arquitectura hexagonal nos daba 
el equilibrio entre dificultad de desarrollo, mantenibilidad y escalabilidad que buscábamos para ClimbIt, permitiéndonos
desacoplar toda la lógica de negocio de factores externos como la base de datos y el frontend, quedando así un sistema 
más modular y fácil de mantener a largo plazo.

En cuanto al stack tecnológico, para el backend optamos por usar Node.js con Express debido a su popularidad, fiabilidad, 
facilidad de uso y su amplio ecosistema de librerías. Elegimos trabajar directamente con JavaScript y en formato ESModules 
y no con TypeScript principalmente por la velocidad de desarrollo y porque Node.js no tiene un soporte nativo para TypeScript.
Para la base de datos, descartamos desde el principio las bases 
de datos no relacionales, ya que la naturaleza de los datos de ClimbIt, con relaciones claras entre escaladores, 
rocódromos y rutas se adaptaba mejor a un modelo relacional. Ambos habíamos trabajado previamente con bases de datos 
como MariaDB y MySQL, pero debido a su rendimiento, escalabilidad y su más que creciente popularidad y ecosistema 
decidimos que para nuestro proyecto, y para las ideas que queremos llegar a implementar en el futuro, PostgreSQL 
era la mejor opción. Para la comunicación entre el backend y la base de datos, optamos por usar un ORM 
(\textit{Object-Relational Mapping}) para facilitar el desarrollo, así como una capa extra de seguridad. Después de evaluar varias
opciones elegimos que Sequelize, debido a su diseño nativo para JavaScript, era la opción más indicada. Otros ORMs como 
TypeORM o Prisma también eran opciones válidas, pero al estar más orientados a TypeScript no podríamos aprovechar al máximo 
sus características, por lo que Sequelize era la mejor opción para nuestro proyecto.

A la hora de servir imágenes al frontend, llegamos a la conclusión de que tener imágenes sin procesar
en el servidor jugaría en nuestra contra a la hora de tener un buen rendimiento, por lo que decidimos generalizar el uso
del formato webp en todas las imágenes que se subieran a la aplicación, exceptuando los mapas de las zonas, 
que se mantendrían en formato svg. Obligar a que los usuarios subieran las imágenes en un único formato generaría 
una gran fricción a la hora de usar la app; la solución más óptima era permitir una amplia gama de formatos para las imágenes
y convertirlas a webp en el servidor antes de guardarlas. Para ello decidimos usar la librería ffmpeg, una de 
las librerías más populares y potentes para el procesamiento de imágenes que ya habíamos utilizado previamente en otro proyecto.

En cuanto a la gestión de sesiones, optamos por JWT (\textit{JSON Web Tokens}) por su 
ligereza y capacidad de escala futura. Los tokens contienen la información necesaria sobre 
permisos e identidad del usuario, eliminando la necesidad de realizar validaciones completas con cada petición 
y mejorando así los tiempos de respuesta y el rendimiento. Además, al estar firmados con una clave secreta, garantizamos 
que los datos no pueden ser manipulados, proporcionando una capa adicional de seguridad sin 
sobrecargar el servidor con sesiones almacenadas.

Para el frontend sabíamos que queríamos una aplicación web, pero que fuera PWA (\textit{Progressive Web App}), para que los usuarios 
de cualquier dispositivo pudieran acceder a ella y descargarla sin la necesidad de tener que crear una aplicación nativa para 
cada plataforma. Desde el principio fuimos reticentes con el uso de frameworks frontend. Después de valorar las ventajas que 
nos ofrecían frente a la curva de aprendizaje y la complejidad añadida que suponían, decidimos optar por una solución más ligera 
y sencilla: utilizar HTML, CSS y JavaScript puro, ayudados de Bootstrap para mantener un diseño consistente. Esto nos permitió 
tener un control total sobre el código y una velocidad inicial de desarrollo mucho mayor, ya que los dos habíamos desarrollado 
previamente con estas tecnologías. Por otro lado, sí nos supuso un reto la implementación de ciertas funcionalidades, en 
concreto el enrutamiento de las vistas o la implementación manual de componentes personalizados, que es algo que los 
frameworks manejan de forma nativa y que nosotros tuvimos que implementar desde cero, pero al final creemos que fue la decisión 
correcta para el proyecto. 


\section{Diseño de la Arquitectura}

\subsection{Backend}

Una vez elegida la arquitectura hexagonal para el backend, se diseñó la estructura del proyecto siguiendo los principios de 
esta arquitectura y aplicando una serie de patrones de diseño que nos permitieran mantener un código limpio, sencillo y modular, 
completamente independiente entre las diferentes capas. La arquitectura hexagonal organiza el código en capas definidas, 
donde cada una tiene una responsabilidad clara. La capa de dominio se encarga de la lógica de negocio y las reglas del sistema, 
donde se definen todos los contratos y validaciones sobre los objetos que se procesan. La capa de infraestructura se encarga
de la comunicación con la base de datos, que nosotros gestionamos a través de Sequelize, así como los modelos que representan
las entidades del sistema, las migraciones y las seeds. La capa de aplicación se encarga de definir la lógica de los use 
cases, orquestando la interacción entre la capa de dominio y la capa de infraestructura. Por último, la capa de presentación, 
que nosotros hemos llamado interfaz, se encarga de la comunicación con el frontend a través de una API REST, definiendo los 
endpoints, los controladores que gestionan las peticiones y respuestas y los middlewares que nos ayudan a la autenticación 
de tokens, formatos de datos y validaciones tanto de datos como de archivos.

{Diagrama de la arquitectura hexagonal del backend con ejemplo de ClimbIt}

A lo largo del desarrollo nos encontramos con diversas situaciones que nos llevaron a aplicar diferentes patrones de diseño que 
nos ayudaron a obtener un código más limpio, modular y fácil de mantener. Los patrones de diseño que hemos utilizado son los siguientes:

\paragraph{Patrón Repository:} Este patrón propuesto por Martin Fowler consiste en crear una interfaz entre la lógica de negocio y la capa
de persistencia, en nuestro caso, nuestra base de datos PostgreSQL a través de Sequelize, de manera que nos permite desacoplar el negocio
de la infraestructura, manteniendo así un código más mantenible en el tiempo. En nuestro caso, como se puede ver en la Figura XX,
la capa de dominio define la interfaz del repositorio en la que se establecen los métodos que se necesitan para realizar las diferentes operaciones
y la capa de infraestructura implementa dicha interfaz con la que hablar con la base de datos. De esta manera, la capa de aplicación
solo se encarga de definir la lógica de los distintos casos de uso.

{Figura del diagrama de patrón repository}

\paragraph{Patrón Adapter:} El patrón Adapter es uno de los patrones clásicos introducidos por Erich Gamma, Richard Helm, Ralph Johnson y John Vlissides, 
también conocidos como los ``Gang of Four'' y se encarga de adaptar distintas interfaces para que puedan trabajar entre ellas. En nuestro caso,
lo hemos utilizado para adaptar distintos errores que pueden surgir en las diferentes capas, como se puede observar en la Figura XX, para 
centralizar el manejo de estos errores en un único lugar, permitiéndonos así mantener una consistencia a la hora de tratar los distintos 
errores que pueden surgir.

{Figura del diagrama de patrón adapter}

\paragraph{Patrón Singleton:} Introducido también por ``The Gang of Four'' en su libro ``Design Patterns: Elements of Reusable Object-Oriented Software'', 
el patrón Singleton sirve para garantizar que una clase tenga únicamente una instancia al mismo tiempo a lo largo de la ejecución de la aplicación. 
En nuestro caso, lo hemos utilizado, por ejemplo, para mantener una única conexión a la base de datos, evitando así posibles problemas y errores relacionados
con múltiples conexiones abiertas, pero también lo hemos utilizado en otros puntos que detallaremos más adelante en la implementación de la aplicación.

\paragraph{Inyección de dependencias e Inversión de control:} La inversión de control es un principio abstracto popularizado por Robert C. Martin, también conocido 
como ``Uncle Bob'', pero ideada inicialmente por Martin Fowler, que viene a referir que los objetos no deben crear sus propias dependencias, sino que le tienen
que venir dadas de fuera. Por otro lado, la inyección de dependencias es una de las técnicas para implementar este principio. En nuestra aplicación hemos
utilizado la inyección de dependencias a través de la inyección por constructor transversalmente en toda la aplicación. El controlador recibe por inyección
los casos de uso que necesita para gestionar las peticiones, los use cases reciben por inyección los repositorios que necesitan para realizar sus operaciones 
y los repositorios reciben por inyección los modelos que necesitan para hablar con la base de datos. De esta manera, hemos conseguido un código completamente desacoplado
alineado con la filosofía de la arquitectura hexagonal.

\subsection{Frontend}

Se describe cómo se organiza la interfaz de usuario y su comunicación con el backend.

\section{Diseño de la Base de Datos}

Como se mencionó anteriormente en la sección de decisiones tecnológicas, la naturaleza de los datos en ClimbIt presentaba 
relaciones claras y bien definidas entre entidades como escaladores, rocódromos, zonas y rutas, lo que nos llevó a 
descartar desde el principio las bases de datos no relacionales en favor de un modelo relacional. La estructura jerárquica 
y las múltiples relaciones muchos-a-muchos entre entidades hacían que un modelo relacional fuese la opción más apropiada 
para garantizar integridad referencial, coherencia de datos y facilidad de consultas complejas.

Elegimos PostgreSQL como sistema gestor de bases de datos relacional (SGBDR) debido a su robustez, rendimiento, escalabilidad 
y su creciente popularidad en el ecosistema de desarrollo. Para facilitar la interacción entre la aplicación y la base de datos, 
utilizamos Sequelize como ORM, que nos permite trabajar con los modelos de datos de forma orientada a objetos sin necesidad de 
escribir SQL directamente, proporcionando además una capa adicional de seguridad contra inyecciones SQL y abstrayendo las 
particularidades específicas de PostgreSQL.

\subsection{Entidades Principales del Sistema}

El modelo de datos de ClimbIt está compuesto por las siguientes entidades principales:

\begin{itemize}
    \item \textbf{Escalador:} Representa a los usuarios del sistema, almacenando información de perfil, credenciales de 
    autenticación y configuración personal.
    
    \item \textbf{Rocódromo:} Entidad que representa cada uno de los lugares donde se practica la escalada, incluyendo información 
    como nombre, ubicación y datos de contacto.
    
    \item \textbf{Zona:} Las diferentes áreas o sectores dentro de un rocódromo, cada una con su propio mapa y características.
    
    \item \textbf{Pista:} Las rutas específicas de escalada dentro de una zona, con información de dificultad, tipo de escalada 
    y estado.
    
    \item \textbf{Ascentso:} Registro de cada vez que un escalador asciende una pista, almacenando información como fecha, 
    valoración y estado del ascenso.
    
    \item \textbf{EscalaDificultad:} Tabla de catálogo que define los niveles de dificultad disponibles en el sistema.
    
    \item \textbf{Amistad:} Gestión de relaciones entre escaladores, permitiendo crear conexiones entre usuarios.
    
    \item \textbf{Suscripción:} Control de suscripciones de escaladores a rocódromos específicos.
    
    \item \textbf{FotosPerfil:} Almacenamiento de imágenes de perfil de escaladores.
\end{itemize}

\subsection{Relaciones Principales del Modelo}

El modelo relacional de ClimbIt está estructurado mediante las siguientes relaciones:

\paragraph{Relaciones Jerárquicas:} La estructura principal del sistema sigue una jerarquía clara:
\begin{itemize}
    \item Un Rocódromo contiene múltiples Zonas (relación 1:N).
    \item Una Zona contiene múltiples Pistas (relación 1:N).
\end{itemize}
Esta jerarquía refleja la estructura física de los rocódromos reales en la vida.

\paragraph{Relaciones Muchos-a-Muchos:} El sistema implementa varias relaciones M:N que requieren tablas de unión:
\begin{itemize}
    \item Escalador-Escalador a través de Amistad, permitiendo que dos usuarios establezcan relaciones bidireccionales.
    \item Escalador-Rocódromo a través de Suscripción, permitiendo que un escalador se suscriba a múltiples rocódromos 
    y que un rocódromo tenga múltiples suscriptores.
    \item Escalador-Pista a través de Ascentso, registrando cada ascenso realizado por cada escalador.
\end{itemize}

\paragraph{Relaciones de Integridad Referencial:} Todas las claves foráneas están configuradas con restricciones 
de integridad referencial adecuadas para mantener la consistencia de datos, evitando operaciones que dejasen la 
base de datos en un estado inconsistente.

\subsection{Normalización y Diseño}

La estructura de la base de datos ha sido diseñada siguiendo los principios de normalización relacional, alcanzando 
la tercera forma normal (3FN) en la mayoría de las entidades. Esto implica:

\begin{itemize}
    \item Eliminación de datos redundantes, evitando la duplicación de información.
    \item Cada atributo depende únicamente de la clave primaria de su entidad.
    \item Separación de conceptos en entidades distintas con relaciones apropiadas entre ellas.
\end{itemize}

Estas prácticas de normalización garantizan que las operaciones de actualización, inserción y eliminación sean eficientes 
y no generen inconsistencias en los datos.

\subsection{Decisiones de Diseño Específicas}

A lo largo del diseño de la base de datos tomamos varias decisiones específicas:

\paragraph{Campos de Estado y Booleanos:} Utilizamos campos booleanos para representar estados como \texttt{Activo} en 
rocódromos, zonas y pistas, permitiendo desactivar entidades sin eliminarlas completamente de la base de datos, 
facilitando así el mantenimiento histórico de datos.

\paragraph{Gestión de URLs y Rutas de Imágenes:} Campos como \texttt{logoURL}, \texttt{imagenURL} y \texttt{mapa} almacenan 
rutas a los recursos multimedia asociados a cada entidad. Las imágenes se almacenan en el servidor tras ser procesadas 
a formato webp mediante ffmpeg, mejorando rendimiento sin perder calidad.

\paragraph{Timestamps de Auditoría:} Cada entidad cuenta con campos \texttt{FechaCreacion} y, en casos relevantes, 
\texttt{FechaRetirada} o \texttt{FechaModificacion} para mantener un registro histórico de cambios en los datos.

\paragraph{Enumeraciones:} Ciertos campos utilizan tipos enumerados como \texttt{Estado} en ascensos (completado, intentando, etc.) 
y \texttt{Tipo} en pistas (búlder, ruta, etc.), proporcionando validación a nivel de base de datos.

\subsection{Migraciones y Versionado}

La estructura de la base de datos se gestiona a través de migraciones controladas por Sequelize, permitiendo:

\begin{itemize}
    \item Versionado de cambios en el esquema de la base de datos.
    \item Posibilidad de revertir cambios si es necesario.
    \item Documentación implícita de la evolución del modelo de datos.
    \item Reproducibilidad del esquema en diferentes entornos (desarrollo, pruebas, producción).
\end{itemize}

Adicionalmente, se utilizan seeds (\textit{seeders}) para poblar la base de datos con datos iniciales como niveles de 
dificultad estándar, tipos de escalada y datos de prueba, facilitando el desarrollo y las pruebas del sistema.

\section{Diseño de la UI y UX}

Se recogen los criterios de interacción visual, flujo de pantallas y experiencia de uso.